\section{Conclusion}
\label{sec:conclusion}

We have presented \GreenFaaS, a carbon-aware scheduling framework for
serverless workloads. The technical core is
Lemma~\ref{lemma:tradeoff}, a closed-form characterization of the
cold-start carbon trade-off, and the recognition
(\S\ref{sec:tradeoff:idle}) that the same dimensionless idle-cost
ratio governs the analogous temporal-shift decision. The resulting
SLA-tiered algorithm has $O(R \cdot N)$ per-invocation cost and no
cross-invocation coordination beyond an arrival-rate estimator.

On real LWA 2020 carbon traces, \GreenFaaS{} reduces operational
carbon by 64--83\% on multi-region topologies with a low-carbon
refuge, by 52--53\% on coal-belt topologies with limited spatial
opportunity, and by 0\% (matching FIFO by design) on single-region
topologies where no carbon-aware opportunity exists.
\GreenFaaS{} matches pure spatial routing within 1.3
percentage points across all real-carbon topologies; it does not
demonstrate strict dominance over Spatial on any real-data
configuration we tested. \GreenFaaS's distinct contribution relative
to Spatial is \emph{topology-robustness} (matching Spatial when
Spatial works, falling back to FIFO byte-for-byte when it does not)
and the principled cold-start trade-off underlying its decisions,
not a larger carbon reduction.

Three open directions stand out for follow-up: the multi-container
generalization of Lemma~\ref{lemma:tradeoff} (\S\ref{sec:discussion}),
a FaaS-adapted comparison against the
\citet{lechowicz2024online} double-threshold algorithm, and a
production-system validation on Knative. The simulator, trace
loaders, scheduler implementations, Caribou port, and manuscript
sources are released as open-source artifacts under the Apache-2.0
license upon publication, to support all three.
